\section{Stability of the FSA-v2 SDE: a Lyapunov analysis}\label{sec:lyap}

Identifiability tells us that the inference problem for $\theta$ is
well-posed; stability tells us that the dynamical system itself is
well-behaved. This section asks two stability questions about the
FSA-v2 SDE in its G1-reparametrized form: where are the equilibria
of its deterministic skeleton, and is the full stochastic process
ergodic on its physiological operating domain? The first question
admits a clean answer in terms of the Stuart--Landau growth rate;
the second is settled affirmatively via a quadratic Lyapunov
function. Detailed Jacobian computations, the Foster--Lyapunov drift
calculation, and bounds on constants are deferred to
Appendix~\ref{app:lyap}.

\subsection{Deterministic skeleton: equilibria and the Stuart--Landau bifurcation}\label{sec:lyap-det}

Set $\sigma_B = \sigma_F = \sigma_A = 0$ in
\eqref{eq:fsa-B}--\eqref{eq:fsa-A} and a constant control
$\Phi(t) \equiv \Phi_c$. The deterministic skeleton has an explicit
$A=0$ branch
\begin{equation}
  B^*_0(\Phi_c) = \frac{\kappa_B \tau_B \Phi_c}{1+\varepsilon_A A_{\mathrm{typ}}},
  \quad
  F^*_0(\Phi_c) = \kappa_F \tau_F (1+\lambda_A A_{\mathrm{typ}})\,\Phi_c,
  \quad
  A^*_0 = 0,
  \label{eq:eq-zero}
\end{equation}
together with an implicit oscillatory branch where $A^2 = \mu(B,F)/\eta$.
Local stability of the $A=0$ branch is determined entirely by the
sign of the Stuart--Landau growth rate at the equilibrium, as the
following result records.

\begin{proposition}[Local stability of the $A=0$ branch]\label{prop:stab-zero}
The deterministic equilibrium \eqref{eq:eq-zero} is locally
asymptotically stable if and only if
$\mu(B^*_0(\Phi_c), F^*_0(\Phi_c)) < 0$.
\end{proposition}

The proof is short: the drift Jacobian at \eqref{eq:eq-zero} turns
out to be upper triangular, with two unconditionally negative
eigenvalues coming from the slow Banister directions and a third
eigenvalue equal to $\mu(B^*_0,F^*_0)$ controlling the autonomic
amplitude. Details are in Appendix~\ref{app:lyap}.

The bifurcation locus $\{\Phi_c \!:\! \mu(B^*_0(\Phi_c),F^*_0(\Phi_c)) = 0\}$
turns out to be a single positive root
$\Phi_c^{\mathrm{crit}} \approx 7.5$ TRIMP-units per day at the
truth values of Table~\ref{tab:fsa-truth}. The picture this paints
is physiologically intuitive: at the canonical baseline
$\Phi_c = 1$, $\mu(B^*_0,F^*_0) \approx 0.13 > 0$, so the $A=0$
equilibrium is unstable and the model predicts that an athlete
training at the default load settles on a non-trivial autonomic
oscillation. Above the critical training load
$\Phi_c \approx 7.5$, the quadratic overtraining-collapse term
$\mu_{FF}(F-F_{\mathrm{typ}})^2$ dominates, $\mu$ falls back through
zero, and the system loses its oscillation: the autonomic amplitude
collapses to zero, consistent with empirical observations that very
large chronic training loads suppress autonomic variability.

\subsection{Stochastic stability and exponential ergodicity}\label{sec:lyap-stoch}

We turn to the full SDE \eqref{eq:fsa-B}--\eqref{eq:fsa-A} on the
operating domain
$\Xcal = [0,1] \times [0,\infty) \times [0,\infty)$. The
state-dependent diffusions vanish at the boundary $\partial\Xcal$,
so $\Xcal$ is positively invariant. We seek a single quantitative
statement that captures the long-time behaviour of the process. The
Foster--Lyapunov framework supplies the right object.

\begin{definition}[Foster--Lyapunov drift]\label{def:foster-lyap}
A function $V:\Xcal \to [0,\infty)$ is said to satisfy a
Foster--Lyapunov drift inequality with rate $\alpha > 0$ and
constant $C \ge 0$ if the infinitesimal generator $\Lcal$ of the
SDE satisfies
\begin{equation}
  \Lcal V(x) \;\le\; -\alpha\,V(x) + C
  \quad\text{for all } x \in \Xcal.
  \label{eq:foster-lyap}
\end{equation}
\end{definition}

Inequality \eqref{eq:foster-lyap} is the standard hypothesis that
delivers exponential ergodicity in the sense of total-variation
convergence to a unique invariant distribution
(Khasminskii, 1980; Meyn and Tweedie, 1993). We exhibit a $V$ that
satisfies it.

\begin{theorem}[Stochastic stability of FSA-v2]\label{thm:lyap}
Under bounded controls $\Phi:\R_+ \to [0,\Phimax]$, the FSA-v2 SDE
\eqref{eq:fsa-B}--\eqref{eq:fsa-A} satisfies the Foster--Lyapunov
drift inequality \eqref{eq:foster-lyap} for the quadratic Lyapunov
function
\begin{equation}
  V(B,F,A) \;=\; \tfrac{1}{2}\bigl(a B^2 + b F^2 + c A^2\bigr),
  \qquad a,b,c > 0,
  \label{eq:V-quadratic}
\end{equation}
with rate
$\alpha = \min\!\bigl(2/\tau_B,\,1/[\tau_F(1+\lambda_A A_{\mathrm{typ}})]\bigr)$
and a finite constant $C$ depending on the truth parameters,
$\Phimax$, and the choice of $a,b,c$.
\end{theorem}

The proof is a routine but tedious computation of $\Lcal V$ component
by component, bounding each term using the boundedness of $B$ on
$[0,1]$ together with two applications of Young's inequality on the
$F$ and $A$ components; we record it in
Appendix~\ref{app:lyap}. The intuitive picture is that the three
dissipative terms in the drift (linear decay of $B$, linear decay
of $F$, cubic damping of $A$) overwhelm the linear growth and noise
contributions outside a bounded region of state space, which is
exactly what the Foster--Lyapunov inequality formalises.

\begin{theorem}[Khasminskii, 1980; Meyn--Tweedie, 1993]\label{thm:ergodic}
Let $X_t = (B_t,F_t,A_t)$ be the FSA-v2 process under a fixed bounded
control $\Phi$, satisfying the conditions of
Theorem~\ref{thm:lyap}. Then $X$ admits a unique invariant
probability distribution $\pi^\Phi$ on $\Xcal$, and there exist
constants $K_0 > 0$ and $\rho \in (0,\alpha)$ such that, for every
$x_0 \in \Xcal$,
\begin{equation}
  \bigl\| \mathrm{Law}(X_t \mid X_0 = x_0) - \pi^\Phi \bigr\|_{\mathrm{TV}}
   \;\le\; K_0\,(1 + V(x_0))\,e^{-\rho t}
  \quad \text{for all } t \ge 0,
  \label{eq:ergodic-bound}
\end{equation}
where $\|\cdot\|_{\mathrm{TV}}$ is the total-variation norm.
\end{theorem}

The exponential rate $\rho$ inherits the slowest dissipative
timescale of the model. With $\tau_B = 42$ days and
$\tau_F^{\mathrm{old}} := \tau_F(1+\lambda_A A_{\mathrm{typ}}) = 7$ days,
the rate is $\alpha = 1/7\,\mathrm{d}^{-1}$ and convergence to
invariance has half-life roughly $7\log 2 \approx 4.85$ days. This
is consistent with the rolling-window choice of one-day windows
and twelve-hour strides in Section~\ref{sec:smc2}: the
SMC${}^2$ filter operates on a fundamental dynamical scale
comparable to the relaxation time of the process being filtered.

\begin{remark}[Limitations of the bound]\label{rem:lyap-limits}
The proof of Theorem~\ref{thm:lyap} in Appendix~\ref{app:lyap} uses
the simplest possible quadratic Lyapunov function and crude global
bounds on the residual factors of the drift. The constant $C$ that
emerges is therefore loose: it scales with $\Phimax^2$ rather than
with the time-average of $\Phi^2$, and ignores the cancellations
that the real dynamics exhibit near the operating point. A sharper
analysis -- perhaps with a piecewise Lyapunov function that
switches between the regimes $\mu < 0$ and $\mu > 0$ identified in
Proposition~\ref{prop:stab-zero} -- would deliver a tighter rate
and a smaller constant. Such a refinement is left for future work.
\end{remark}
